
\examquestion{2024}{2}{Zeeman-transition absorption as a magnetic-field sensor}{3}
The problem supplies the steady-state Bloch inversion
\[
\langle\sigma_z\rangle_{\rm ss}
=-\frac{\gamma^2+4\Delta(B)^2}{
\gamma^2+4\Delta(B)^2+8\Omega^2},\qquad
\Delta(B)=\kappa B.
\]
Here \(\Omega\) and \(\gamma\) are \emph{frequencies};
the Hamiltonian's drive term must have energy units,
i.e.\ effectively \(\hbar\Omega\sigma_x\), if the
printed shorthand is missing \(\hbar\).

\subsection*{(a) Probability of the ground state}
Since \(\sigma_z=\ket e\bra e-\ket g\bra g\),
\(P_e-P_g=\langle\sigma_z\rangle\) while \(P_e+P_g=1\).
Therefore \(P_g=(1-\langle\sigma_z\rangle)/2\).
Define \(D=\gamma^2+4\kappa^2 B^2+8\Omega^2\), giving
\[
P_g=\frac12\left(1+\frac{D-8\Omega^2}{D}\right)
=\boxed{1-\frac{4\Omega^2}{
\gamma^2+4\kappa^2B^2+8\Omega^2}}.
\]
On resonance \(B=0\), \(P_g\) is smallest; well away
from resonance \(P_g\to1\).

\subsection*{(b) What is measured?}
The AC absorption signal is assumed proportional to \(P_g\).
Calibrating the signal versus \(B\) gives a Lorentzian-shaped
dip. For unknown perturbations near a known bias \(B_0\),
\[
P_g(B_0+\delta B)\simeq P_g(B_0)
+\left.\frac{\diff P_g}{\diff B}\right|_{B_0}\delta B.
\]
Measuring the absorption change then estimates
\(\delta B\simeq \delta P_g/P'_g(B_0)\).
The signal is \emph{even} in \(B\); its sign alone cannot
distinguish \(+B\) from \(-B\). A known bias removes that ambiguity.

\subsection*{(c) Determine the optimal linear operating point}
Put \(C=\gamma^2+8\Omega^2\). Differentiate:
\[
\frac{\diff P_g}{\diff B}
=\frac{32\Omega^2\kappa^2B}
{(C+4\kappa^2B^2)^2}.
\]
The sensitivity magnitude is maximised where
the second derivative is zero:
\[
\frac{\diff^2 P_g}{\diff B^2}
=32\Omega^2\kappa^2
\frac{C-12\kappa^2B^2}{(C+4\kappa^2 B^2)^3}=0.
\]
Thus \(12\kappa^2B_0^2=C\), so
\[
\boxed{B_0=\pm\frac{\sqrt{\gamma^2+8\Omega^2}}
{2\sqrt3\,|\kappa|}.}
\]
\begin{cautionbox}
\textbf{Dimensional check.} The numerator has units of
frequency, so the expression for \(B_0\) must include
division by \(\kappa\) (frequency per magnetic field).
If the source print or text extraction omits this factor,
the formula as printed cannot be dimensionally correct.
The result above follows unambiguously from the supplied
\(\Delta(B)=\kappa B\).
\end{cautionbox}
\solcheck{The slope vanishes at \(B=0\) and far off
resonance. The inflection points at \(\pm B_0\) produce
maximum local linear response.}
